% =====================================================================
% 5.4 Image provenance and watermark robustness (optional lab)
%
% Source: ../aisb/5.4-watermarking/section4_solution.py
% Figures (from the Tree-Ring arXiv source, Wen et al. 2023):
%   treering_teaser.pdf        Fig. 1 pipeline
%   treering_attack_*.png      appendix "Attacked images"; parameters in
%                              the captions are the paper's, not the lab's
% The lab records no detection rates. The cluster values (about 45 vs
% about 20) are quoted from the lab prose; robustness outcomes are the
% lab's own summary bullets.
% =====================================================================

\section{5.4: Watermarking (optional)}

\begin{frame}
\frametitle{5.4: switch sides}
\begin{center}
\includegraphics[width=0.78\textwidth]{treering_teaser}
\figsrc{Tree-Ring watermarking. Figure from \textcite{wen2023treering}.}
\end{center}
\begin{itemize}
    \item Question: was this image made by \textit{our} model?
    \item The defender now needs a signal that survives an attacker who edits
    the image.
\end{itemize}
\end{frame}

\begin{frame}
\frametitle{Where the key goes: the starting noise}
A diffusion model turns random noise $x_T$ into an image. Choose that noise
carefully instead of at random.
\begin{block}{Embedding}
Take the 2-D Fourier transform $\Fou(x_T)$. Inside a small disk around the
centre (the lowest frequencies), overwrite the coefficients with a secret
pattern of concentric rings. Transform back, then generate as normal.
\end{block}
\begin{itemize}
    \item Low frequencies are the coarse, large-scale structure. Blur, JPEG
    and resizing mostly throw away high frequencies. Key remains undamaged.
    \item Rings are rotation invariant: rotating an image rotates its spectrum,
    and a ring rotated is the same ring.
\end{itemize}
\end{frame}

\begin{frame}
\frametitle{Detection: run the model backwards}
\begin{itemize}
    \item Can run diffusion backwards:
    image $\to$ estimate of the starting noise $x_T$.
    \item Take $\Fou$ of the recovered noise and measure the distance to the key
    inside the disk.
\end{itemize}
\begin{block}{Hypothesis test}
$$
\text{watermarked} \iff \big|\, \Fou(\hat x_T) - \text{key} \,\big| \;<\; \tau
$$
\end{block}
\begin{itemize}
    \item In the lab, clean images and wrong keys score around 45, watermarked
    images around 20. The threshold $\tau$ goes between the two clusters.
\end{itemize}
\end{frame}

\begin{frame}
\frametitle{What survives}
\begin{center}
\setlength{\tabcolsep}{2pt}
\begin{tabular}{cccc}
\includegraphics[width=0.2\textwidth]{treering_attack_jpeg} &
\includegraphics[width=0.2\textwidth]{treering_attack_blur} &
\includegraphics[width=0.2\textwidth]{treering_attack_rotate} &
\includegraphics[width=0.2\textwidth]{treering_attack_crop} \\
\scriptsize JPEG 25\% & \scriptsize blur & \scriptsize rotate $75^\circ$ &
\scriptsize crop 75\% + rescale
\end{tabular}
\figsrc{Attacked images from \textcite{wen2023treering}. The lab uses milder settings.}
\end{center}
\begin{itemize}
    \item Blur, JPEG, resizing and noise: survive, because they mostly remove
    high frequencies.
    \item Rotation: survives, because of the rings.
    \item Crop and rescale breaks it. Rescaling stretches the spectrum, so ring
    10 lands where ring 7 or 8 should be. The paper reports the same weakness.
\end{itemize}
\end{frame}

\begin{frame}
\frametitle{The real weakness: need unique keys per image}
\begin{itemize}
    \item Every watermarked image carries the same pattern. Invert ten of them
    and average: the image content average out, and we can recover the key.
    \item Write the recovered key into new noise and generate. The forgeries
    pass the detector as well as genuine images do.
    enough.
    \item Later schemes vary the pattern per image, e.g.\ Gaussian Shading
    \parencite{yang2024gaussian}.
\end{itemize}
\vspace{0.2cm}
\begin{block}{Two more limits}
Detection requires model weights, so users cannot test for AI by running something
locally. Requires API.
\end{block}
\end{frame}
